%%%%%%%%%%%%%%%%%%%%%%%%%%%%%%%%%%%%%%%%%%%%%%%%%%%%%%%%%%%%%%%%%%%%%%%%
%% Renato Mota Costa --- Engenheiro de Dados (versão PT-PT)
%% Construído com AltaCV (v1.7.4) de LianTze Lim.
%%
%% Compilar com:  pdflatex renato-cv-pt.tex
%% (ou)           xelatex -shell-escape -output-driver="xdvipdfmx -z 0" renato-cv-pt.tex
%%
%% NOTA: o babel é carregado DEPOIS do hyperref (que vem da própria classe).
%% Normalmente funciona sem problemas, mas se aparecer algum erro relacionado
%% a bookmarks ou acentos, basta comentar a linha do babel --- perde-se apenas
%% a hifenização portuguesa.
%%%%%%%%%%%%%%%%%%%%%%%%%%%%%%%%%%%%%%%%%%%%%%%%%%%%%%%%%%%%%%%%%%%%%%%%

\documentclass[10pt,a4paper,ragged2e,withhyper]{altacv}

%% Layout da página
\geometry{left=1.25cm,right=1.25cm,top=1.4cm,bottom=1.4cm,columnsep=1.1cm}

%% Colunas paralelas que quebram entre páginas
\usepackage{paracol}

%% Hifenização portuguesa (ver NOTA acima)
\usepackage[portuguese]{babel}

%% Tipos de letra
\iftutex
  \setmainfont{Roboto Slab}
  \setsansfont{Lato}
  \renewcommand{\familydefault}{\sfdefault}
\else
  \usepackage[rm]{roboto}
  \usepackage[defaultsans]{lato}
  \renewcommand{\familydefault}{\sfdefault}
\fi

%%%%%%%%%%%%%%%%%%%%%%%%%%%%%%%%%%%%%%%%%%%%%%%%%
%% Paleta de cores --- azul-marinho + azul oceano
%%%%%%%%%%%%%%%%%%%%%%%%%%%%%%%%%%%%%%%%%%%%%%%%%
\definecolor{DeepNavy}{HTML}{0B2545}
\definecolor{OceanBlue}{HTML}{0E6BA8}
\definecolor{SteelBlue}{HTML}{3A7CA5}
\definecolor{SlateDark}{HTML}{2E3440}
\definecolor{BodyGrey}{HTML}{4C5565}

\colorlet{name}{DeepNavy}
\colorlet{tagline}{OceanBlue}
\colorlet{heading}{DeepNavy}
\colorlet{headingrule}{SteelBlue}
\colorlet{subheading}{OceanBlue}
\colorlet{accent}{OceanBlue}
\colorlet{emphasis}{SlateDark}
\colorlet{body}{BodyGrey}

%% Tipos de letra dos elementos estruturais
\renewcommand{\namefont}{\Huge\rmfamily\bfseries}
\renewcommand{\taglinefont}{\large\bfseries}
\renewcommand{\personalinfofont}{\footnotesize\bfseries}
\renewcommand{\cvsectionfont}{\Large\rmfamily\bfseries}
\renewcommand{\cvsubsectionfont}{\normalsize\bfseries}

%% Marcadores
\renewcommand{\cvItemMarker}{{\small\textbullet}}
\renewcommand{\cvRatingMarker}{\faCircle}
\renewcommand{\cvDateMarker}{\faCalendar[regular]}
\renewcommand{\cvLocationMarker}{\faMapMarker}

%% Texto alternativo dos ícones de data/local (ajuda em sistemas ATS:
%% quem copia o PDF em texto simples lê "Data:" e "Local:")
\renewcommand{\locationname}{Local}
\renewcommand{\datename}{Data}

%%%%%%%%%%%%%%%%%%%%%%%%%%%%%%%%%%%%%%%%%%%%%%%%%
%% Auxiliar: uma subfunção dentro do bloco de um só empregador
%% #1 = área / divisão     #2 = período
%%%%%%%%%%%%%%%%%%%%%%%%%%%%%%%%%%%%%%%%%%%%%%%%%
\newcommand{\cvsubevent}[2]{%
  \smallskip
  {\color{emphasis}\normalsize\bfseries #1\hfill
   \BeginAccSupp{method=pdfstringdef,ActualText={\datename:}}%
   {\small\color{accent}\cvDateMarker}\EndAccSupp{}%
   {\small\color{accent}~#2}\par}%
  \smallskip
}

%% Listas mais compactas
\setlist[itemize]{label=\cvItemMarker,leftmargin=1.2em,itemsep=0.25em,topsep=0.3em,parsep=0pt}

%%%%%%%%%%%%%%%%%%%%%%%%%%%%%%%%%%%%%%%%%%%%%%%%%
%% Dados do cabeçalho
%%%%%%%%%%%%%%%%%%%%%%%%%%%%%%%%%%%%%%%%%%%%%%%%%
\name{Renato Mota Costa}
\tagline{Engenheiro de Dados \textbullet\ AWS \textbullet\ PySpark \textbullet\ Python \textbullet\ SQL\\[0.35em]
  {\normalsize\mdseries\color{body}Mudança para Portugal --- autorizado a trabalhar (visto D2) \textbullet\ disponibilidade imediata}}

\personalinfo{%
  \email{renato.mota.costa@gmail.com}
  \phone{+55 (11) 96621-3126}
  \linkedin{renato-mota-31890a227}
  \github{Rmcstt}
}

%% Descomentar para incluir fotografia (opção "normalphoto" para foto quadrada)
% \photoR{2.5cm}{ficheiro-da-foto}

\begin{document}
\makecvheader

%% Coluna esquerda 62%, coluna direita 38%
\columnratio{0.62}
\begin{paracol}{2}

%%%%%%%%%%%%%%%%%%%%%%%%%%%%%%%%%%%%%%%%%%%%%%%%%
%% COLUNA ESQUERDA
%%%%%%%%%%%%%%%%%%%%%%%%%%%%%%%%%%%%%%%%%%%%%%%%%

\cvsection{Resumo Profissional}

Engenheiro de Dados com 4 anos de experiência na conceção, modernização e
automatização de \textit{pipelines} de dados em larga escala, em ambientes
bancários \textit{cloud-first}. Especializado em AWS (Glue, Lambda, S3, Athena,
Step Functions), PySpark e boas práticas de ETL/ELT, com forte foco em
desempenho, otimização de custos e governação de dados. Pós-graduado em
Engenharia de \textit{Machine Learning}, certificado AWS Cloud Practitioner e
em preparação para a certificação Databricks.

\cvsection{Experiência Profissional}

\cvevent{Engenheiro de Dados}{Itaú Unibanco}{Set 2022 -- Presente}{São Paulo, Brasil}

\cvsubevent{Divisão de Crédito Imobiliário}{Mar 2026 -- Presente}
\begin{itemize}
  \item Aumentou a cobertura de governação de metadados em
        \textbf{75 pontos percentuais} em mais de \textbf{500 tabelas}, em
        menos de um mês.
  \item Lidera a modernização dos \textit{pipelines} de dados em \textit{cloud}
        que sustentam a análise e a concessão de crédito.
  \item Construiu \textit{pipelines} ELT com AWS Glue, PySpark, Lambda, Athena
        e S3; alargou a automatização, a monitorização de custos e as
        verificações de qualidade de dados.
\end{itemize}

\divider

\cvsubevent{Divisão de Crédito Automóvel}{Mar 2025 -- Mar 2026}
\begin{itemize}
  \item Reduziu o tempo de processamento de ETL através da otimização de
        PySpark e da afinação do AWS Glue.
  \item Geriu conjuntos de dados estruturados e semiestruturados em larga
        escala no S3, para \textit{scoring} de crédito e concessão.
\end{itemize}

\divider

\cvsubevent{Divisão de Consórcios}{Set 2022 -- Mar 2025}
\begin{itemize}
  \item Concebeu e implementou \textit{pipelines} de ETL/ELT escaláveis para
        operações de grande volume.
  \item Migrou fluxos de trabalho \textit{legacy} para a AWS, melhorando o
        desempenho e reduzindo custos.
  \item Automatizou fluxos de ingestão com Glue, Lambda, Step Functions e
        Athena.
\end{itemize}

%%%%%%%%%%%%%%%%%%%%%%%%%%%%%%%%%%%%%%%%%%%%%%%%%
%% COLUNA DIREITA
%%%%%%%%%%%%%%%%%%%%%%%%%%%%%%%%%%%%%%%%%%%%%%%%%
\switchcolumn

\cvsection{Competências Técnicas}

\cvsubsection{Cloud e Plataforma de Dados}
\cvtag{AWS Glue}
\cvtag{Lambda}\\
\cvtag{S3}
\cvtag{Athena}\\
\cvtag{Step Functions}
\cvtag{IAM}\\
\cvtag{Databricks \textit{(em curso)}}

\divider\smallskip

\cvsubsection{Linguagens e Processamento}
\cvtag{Python}
\cvtag{PySpark}\\
\cvtag{Spark}
\cvtag{SQL}
\cvtag{MySQL}

\divider\smallskip

\cvsubsection{Práticas e Ferramentas}
\cvtag{ETL / ELT}
\cvtag{Modelação de Dados}\\
\cvtag{Governação de Dados}\\
\cvtag{Automatização}
\cvtag{Git / GitHub}

\cvsection{Formação}

\cvevent{Pós-Graduação em Engenharia de \textit{Machine Learning}}{FIAP}{Concluída em Jan 2025}{}

\divider

\cvevent{Licenciatura em Análise e Desenvolvimento de Sistemas}{Universidade Anhembi Morumbi}{Currículo concluído}{}

\divider

\cvevent{Licenciatura em Educação Física}{}{Concluída em 2020}{}

\cvsection{Certificações}

%% Se o seu fontawesome5 for antigo e \faAws não existir, troque por \faCloud
\cvachievement{\faAws}{AWS Certified Cloud Practitioner}{Amazon Web Services}

\divider

\cvachievement{\faCertificate}{Certificação Databricks}{Em curso}

\cvsection{Idiomas}

\cvskill{Português --- Nativo}{5}
\divider
\cvskill{Inglês --- B2}{3.5}

\end{paracol}
\end{document}
